\section{Case: The Quant System}

% Page 37
\begin{frame}{Roadmap}
\tableofcontents[currentsection,sectionstyle=show/shaded,subsectionstyle=show/show/hide]
\end{frame}

\subsection{Build}

% Page 38
\begin{frame}{Build: Package and Entry Points}
\begin{center}
\lstinputlisting[style=filetext,basicstyle=\ttfamily\footnotesize,breakatwhitespace=false]{figures/qt_layout.txt}

\end{center}
\begin{itemize}
\item Use \texttt{/qt:start} for research and \texttt{/qt:evo} for improvements.
\item Install with \texttt{/plugins install <qt>}. After source edits, reinstall and run \texttt{/reload}.
\item Keep installed code separate from case state and prompt versions in the workspace's \texttt{.qt/} directory.
\end{itemize}
\end{frame}

% Page 39
\begin{frame}{Build: Leads and Specialists}
\begin{center}
\lstinputlisting[style=filetext,basicstyle=\ttfamily\footnotesize,breakatwhitespace=false]{figures/qt_agent.txt}

\end{center}
\begin{itemize}
\item Define the lead's tools and specialist allowlist in its role card.
\item Inspect actual host dispatch and response records during a live run.
\item Report files alone do not prove that separate specialist contexts ran.
\end{itemize}
\end{frame}

% Page 40
\begin{frame}{Build: Point-in-Time Inputs}
\begin{center}
\small\begin{tabular}{@{}P{3.0cm}P{4.6cm}P{4.9cm}@{}}
\toprule
\textbf{Time} & \textbf{Record} & \textbf{Meaning} \\
\midrule
2024-03-01 & v1: value = 5.0 & Initial release available \\
\rowcolor{gray!10}
2024-05-01 & v2: value = 4.7 & Restated value released \\
April decision & Use v1 = 5.0 & Filter by release date before version \\
\bottomrule
\end{tabular}

\end{center}
\begin{center}\illustrative\end{center}
\begin{itemize}
\item Select the version publicly available at the decision date.
\item A signal calculated from the closing price can only inform a later trade.
\item These are the QT data contract's checks; the harness determines where they run and what happens if they fail.
\end{itemize}
\end{frame}

% Page 41
\begin{frame}{Build: Verified Numerical Tools}
\begin{center}
\small\begin{tabular}{@{}P{3.1cm}P{4.5cm}P{4.9cm}@{}}
\toprule
\textbf{Part} & \textbf{Definition} & \textbf{Implementation} \\
\midrule
Score & $s_{i,t}=P_{i,t}/P_{i,t-20}-1$ & 21 valid closing prices \\
\rowcolor{gray!10}
Example & $12/10-1=20\%$ & Known-answer unit test \\
Rank & Descending score; stable tie-break & \texttt{baseline/strategy.py} \\
\bottomrule
\end{tabular}

\end{center}
\begin{center}\illustrative\end{center}
\lstinputlisting[style=filetext,basicstyle=\ttfamily\footnotesize]{figures/ca_momentum_code.txt}

\begin{center}\scriptsize Source: QT \texttt{baseline/strategy.py}; input validation runs before this function.\end{center}
\begin{itemize}
\item Check dates and price availability before calling the calculation.
\item Test the known answer, verify stable ranking, and state the earliest permitted execution time.
\end{itemize}
\end{frame}

% Page 42
\begin{frame}{Build: The Orchestrator and Gates}
\begingroup
\setstretch{1.0}
\lstinputlisting[style=filetext,basicstyle=\ttfamily\footnotesize]{figures/ca_controller_code.txt}
\endgroup

\begin{center}\scriptsize Source: QT \texttt{orchestrator.py}, excerpt from \texttt{run\_cycle}.\end{center}
\begin{center}
\renewcommand{\arraystretch}{1.0}
\small\begin{tabular}{@{}P{4.0cm}P{8.5cm}@{}}
\toprule
\textbf{Barrier} & \textbf{Evidence Required} \\
\midrule
Generator & All assigned specialist reports \\
\rowcolor{gray!10}
Numerical & Executed candidate and passed gates \\
Critic & Complete review and carried findings \\
\rowcolor{gray!10}
Termination & Acceptance or explicit stop conditions \\
\bottomrule
\end{tabular}

\end{center}
\begin{itemize}
\item Require three matching Generator reports and three fresh Critic reviews.
\item Run \texttt{python3 <qt>/scripts/check.py} to check the controller rules.
\end{itemize}
\end{frame}

% Page 43
\begin{frame}{Build: Register and Inspect a Hook}
\begin{center}
\lstinputlisting[style=filetext,basicstyle=\ttfamily\footnotesize]{figures/qt_hook_manifest.txt}

\end{center}
\begin{itemize}
\item QT's optional Bash logger records event/call IDs and deduplicates repeated deliveries.
\item It neither assesses candidates nor records every specialist action.
\item Kimi logs hook errors and timeouts, then continues. Required checks remain in the controller.
\end{itemize}
\begin{center}\scriptsize\href{https://www.kimi.com/code/docs/en/kimi-code-cli/customization/hooks.html}{Source: Kimi hook behavior}\end{center}
\end{frame}

% Page 44
\begin{frame}{Build: Run One Research Case}
\begin{center}
\lstinputlisting[style=prompt,basicstyle=\ttfamily\footnotesize,breakatwhitespace=false]{figures/qt_start.txt}

\end{center}
\begin{itemize}
\item Follow the \texttt{request\_id}, specialist reports, numerical outputs, and routing decision.
\item Label scripted rehearsals. Show host dispatch records when claiming a live team ran.
\item The final evidence must support acceptance or explain why the case stops.
\item The fixture tests momentum candidates; only \texttt{lookback} and \texttt{top\_n} vary.
\end{itemize}
\end{frame}

% Page 45
\begin{frame}{Build: Repair a Rejected Submission\handsup}
\begin{center}
\small\begin{tabular}{@{}P{4.0cm}P{8.5cm}@{}}
\toprule
\textbf{Submission} & \textbf{Expected Controller Response} \\
\midrule
One specialist missing & Reject; keep the current role and consume an attempt \\
\rowcolor{gray!10}
Stale request identifier & Same rejection and retry rule \\
Second invalid submission & \texttt{BLOCKED}; no advancement \\
\rowcolor{gray!10}
Complete matching reports & Check evidence before advancing \\
\bottomrule
\end{tabular}

\end{center}
\begin{itemize}
\item Question: Remove one required report in a disposable case. What must be repaired before the controller can advance?
\item Check matching reports, artifact hashes, and required conditions. Keep sealed evidence intact.
\end{itemize}
\end{frame}

% Page 46
\begin{frame}[t]{Build: Replay, Audit, and Resume\handsup}
\setlength{\topsep}{0pt}
\begin{center}
\small\begin{tabular}{@{}P{3.0cm}P{9.5cm}@{}}
\toprule
\textbf{Operation} & \textbf{Evidence or Rule} \\
\midrule
Reconstruct & Inputs/code; open/repaired findings; stop rule; role card/version; approval evidence, if retained \\
\rowcolor{gray!10}
Recalculate & Freeze inputs, code and configuration; compare numerical results \\
Fresh model run & New case; retain its own decisions, tool paths and outputs \\
\rowcolor{gray!10}
Resume or reuse & Verify request ID and hashes; resume pending work; a completed case returns its saved result \\
\bottomrule
\end{tabular}

\end{center}
\begin{itemize}
\item Inspect \texttt{state.json}, the route trace, and saved outputs.
\item Start a new case if code, the data snapshot or role instructions change, or for a fresh evaluation.
\item Question: Identify this case's inputs, repairs, stopping rule, prompt version, and approver.
\end{itemize}
\end{frame}

% Page 47
\begin{frame}{Build: Production Data and Adapters}
\begin{center}
\small\begin{tabular}{@{}P{3.0cm}P{9.5cm}@{}}
\toprule
\textbf{Component} & \textbf{Production Requirement} \\
\midrule
Inputs & Licensed, dated market and financial data \\
\rowcolor{gray!10}
Universe & Historical index membership, suspensions and delistings \\
Execution & Lot sizes, price limits, T+1 rules and fees for the relevant board and date \\
\rowcolor{gray!10}
Operations & Host isolation, named responsible people and an independent test of each adapter \\
\bottomrule
\end{tabular}

\end{center}
\begin{itemize}
\item The fixture uses five fictional stocks and synthetic weekday bars, not an exchange calendar or current feed.
\item Complete corporate-action, portfolio, and broker adapters remain to be built.
\end{itemize}
\end{frame}

% Page 48
\begin{frame}{Build: Research Validation}
\begin{center}
\resizebox{.94\textwidth}{!}{\begin{tikzpicture}[
  box/.style={draw=blue,fill=blue!5,rounded corners=2pt,text width=3.3cm,minimum height=1.5cm,align=center,font=\small,inner sep=5pt},
  link/.style={-{Stealth[length=5pt]},thick,draw=black!60}]
\node[box] (n0) at (1.7,0) {\textbf{Baseline input}\\freeze data through $t$};
\node[box] (n1) at (6.6,0) {\textbf{Perturbation}\\alter prices after $t$};
\draw[link] (n0) -- (n1);
\node[box,draw=forestgreen,fill=forestgreen!6] (n2) at (11.5,0) {\textbf{No-leakage check}\\signals $\le t$ unchanged};
\draw[link] (n1) -- (n2);
\end{tikzpicture}
}
\end{center}
\begin{itemize}
\item Test future-price changes, restatements, higher costs, and performance in successive out-of-sample periods.
\item Passing these checks does not establish that the strategy will earn excess returns.
\end{itemize}
\end{frame}

% Page 49
\begin{frame}{Build: Capital Authority and Reconciliation}
\begin{center}
\resizebox{.94\textwidth}{!}{\begin{tikzpicture}[b/.style={draw=blue,fill=blue!6,rounded corners,text width=3.4cm,minimum height=1.0cm,align=center,font=\small},a/.style={-{Stealth},thick}]
\node[b] (n0) {\textbf{DRAFT}\\versioned sheet};
\node[b,right=0.5cm of n0] (n1) {\textbf{REVIEWED}\\human decision};
\draw[a] (n0)--(n1);
\node[b,right=0.5cm of n1] (n2) {\textbf{RECONCILED}\\actual fill import};
\draw[a] (n1)--(n2);
\end{tikzpicture}
}
\end{center}
\begin{itemize}
\item Check portfolio limits, current positions, approval expiry, and actual fills.
\item Research acceptance does not authorize orders. A named person remains responsible for capital use.
\item QT has no broker connection.
\end{itemize}
\end{frame}

\subsection{Evolve}

% Page 50
\begin{frame}{Evolve: Capture the Failure}
\begin{center}
\small\begin{tabular}{@{}P{4.0cm}P{8.5cm}@{}}
\toprule
\textbf{Capture} & \textbf{Expected Record} \\
\midrule
Completed case & Original request and failure artifact \\
\rowcolor{gray!10}
Parent prompt & Exact version before the change \\
Expected behavior & Observation precedes execution \\
\rowcolor{gray!10}
Case split & Development and untouched final cases \\
\bottomrule
\end{tabular}

\end{center}
\lstinputlisting[style=shell,basicstyle=\ttfamily\footnotesize]{figures/ca_capture_command.txt}

\begin{itemize}
\item \texttt{/qt:evo} starts the improvement workflow; \texttt{evo-prepare} links the change request to its case and parent prompt version.
\item Define the timing behavior before testing. Use development cases for feedback and reserve separate final-evaluation cases.
\end{itemize}
\end{frame}

% Page 51
\begin{frame}{Evolve: Propose One Prompt Change\handsup}
\begin{center}
\lstinputlisting[style=diff,basicstyle=\ttfamily\small]{figures/timing_prompt_diff.txt}

\end{center}
\begin{center}\illustrative\end{center}
\begin{itemize}
\item Question: Which instruction addresses the timing omission, and what evidence would show that it worked?
\item This QT interface changes the Generator prompt. Strategy parameters, risk limits, acceptance rules, and controller code remain outside its scope.
\end{itemize}
\end{frame}

% Page 52
\begin{frame}{Evolve: Compare Outputs with promptfoo\handsup}
\begin{center}
\small\begin{tabular}{@{}P{4.0cm}P{8.5cm}@{}}
\toprule
\textbf{Prepared Comparison} & \textbf{Inspect} \\
\midrule
Baseline and candidate prompts & Same cases and model settings \\
\rowcolor{gray!10}
Python assertions & Time ordering; evidence references \\
Output matrix & Each failure and its retained output \\
\bottomrule
\end{tabular}

\end{center}
\begin{center}\illustrative\end{center}
\begin{itemize}
\item Trading is allowed from 30 Sep 09:30 (UTC+8). Only \texttt{PRICE\_01} is retained.
\item Run the companion checks, or load \texttt{promptfooconfig.json} in promptfoo, to compare both outputs.
\item Question: Which output passes both checks?
\end{itemize}
\begin{center}\scriptsize\href{https://www.promptfoo.dev/docs/configuration/guide/}{Source: promptfoo configuration and assertions; external example}\end{center}
\end{frame}

% Page 53
\begin{frame}{Evolve: Generate Candidates with DSPy/GEPA}
\begin{center}
\resizebox{.94\textwidth}{!}{\begin{tikzpicture}[
  box/.style={draw=blue,fill=blue!5,rounded corners=2pt,text width=3.3cm,minimum height=1.5cm,align=center,font=\small,inner sep=5pt},
  link/.style={-{Stealth[length=5pt]},thick,draw=black!60}]
\node[box] (n0) at (1.7,0) {\textbf{Development cases}\\fixed metric};
\node[box] (n1) at (6.6,0) {\textbf{Run + inspect}\\failure feedback};
\draw[link] (n0) -- (n1);
\node[box,draw=forestgreen,fill=forestgreen!6] (n2) at (11.5,0) {\textbf{Propose instruction}\\evaluate candidate};
\draw[link] (n1) -- (n2);
\end{tikzpicture}
}
\end{center}
\begin{itemize}
\item GEPA uses a callable program, development cases, a metric, feedback, and a search budget to propose prompt changes.
\item Test candidates on held-out cases. A score does not authorize adoption.
\end{itemize}
\begin{center}\scriptsize\href{https://github.com/stanfordnlp/dspy/blob/main/docs/docs/diving-deeper/gepa-in-depth.md}{Source: DSPy GEPA; external extension, not integrated into QT}\end{center}
\end{frame}

% Page 54
\begin{frame}{Evolve: Judge the Evidence\handsup}
\begin{center}
\small\begin{tabular}{@{}P{4.0cm}P{8.5cm}@{}}
\toprule
\textbf{Record for Both Versions} & \textbf{Decision Use} \\
\midrule
Fixed-case outputs and failures & Inspect why checks pass or fail \\
\rowcolor{gray!10}
Repeat runs and resource use & Assess variation and operating cost \\
Untouched final cases & Check generalization after selection \\
\rowcolor{gray!10}
Critical control failures & Block adoption regardless of score \\
\bottomrule
\end{tabular}

\end{center}
\begin{center}\illustrative\end{center}
\begin{itemize}
\item The candidate fixes the timing but cites a missing source, \texttt{MISSING\_02}.
\item Question: What must be repaired and independently checked before adopting it?
\end{itemize}
\end{frame}

% Page 55
\begin{frame}{Evolve: Record the Tested Version}
\begin{center}
\small\begin{tabular}{@{}P{4.0cm}P{8.5cm}@{}}
\toprule
\textbf{Registry} & \textbf{What to Retain} \\
\midrule
QT local registry & Prompt versions and active pointer \\
\rowcolor{gray!10}
MLflow counterpart & Prompt versions, evaluations, traces \\
Run configuration & Model, tools, policy, data versions \\
\bottomrule
\end{tabular}

\end{center}
\begin{itemize}
\item QT stores immutable prompt versions and the active pointer in \texttt{.qt/}.
\item A registry alias selects a version; approval needs its own record.
\item MLflow is an external counterpart, not a bundled QT integration.
\end{itemize}
\begin{center}\scriptsize\href{https://www.mlflow.org/docs/latest/genai/prompt-registry/}{Source: MLflow Prompt Registry}\end{center}
\end{frame}

% Page 56
\begin{frame}{Evolve: Approve or Reject Adoption\handsup}
\begin{center}
\resizebox{.94\textwidth}{!}{\begin{tikzpicture}[
  box/.style={draw=blue,fill=blue!5,rounded corners=2pt,text width=3.3cm,minimum height=1.5cm,align=center,font=\small,inner sep=5pt},
  link/.style={-{Stealth[length=5pt]},thick,draw=black!60}]
\node[box] (n0) at (1.7,0) {evo-evaluate\\tests + review};
\node[box] (n1) at (6.6,0) {AWAITING\_APPROVAL\\inspect the evidence};
\draw[link] (n0) -- (n1);
\node[box,draw=forestgreen,fill=forestgreen!6] (n2) at (11.5,0) {evo-adopt\\explicit approval};
\draw[link] (n1) -- (n2);
\end{tikzpicture}
}
\end{center}
\begin{itemize}
\item QT checks code regressions and review fields. A person still needs to assess the comparison evidence.
\item Question: Would you adopt this version? Explain your decision using the saved results.
\end{itemize}
\end{frame}

% Page 57
\begin{frame}{Evolve: Roll Back and Monitor}
\begin{center}
\resizebox{.94\textwidth}{!}{\begin{tikzpicture}[
  box/.style={draw=blue,fill=blue!5,rounded corners=2pt,text width=3.3cm,minimum height=1.5cm,align=center,font=\small,inner sep=5pt},
  link/.style={-{Stealth[length=5pt]},thick,draw=black!60}]
\node[box] (n0) at (1.7,0) {\textbf{Authorize rollback}\\retain both versions};
\node[box] (n1) at (6.6,0) {\textbf{Restore pointer}\\prior prompt active};
\draw[link] (n0) -- (n1);
\node[box,draw=forestgreen,fill=forestgreen!6] (n2) at (11.5,0) {\textbf{Verify behavior}\\run a new case};
\draw[link] (n1) -- (n2);
\end{tikzpicture}
}
\end{center}
\begin{itemize}
\item After approval, use \texttt{evo-rollback} to restore the previous pointer. Keep both versions and their evidence.
\item Existing cases retain their original configuration; they do not switch versions mid-run.
\item Verify the restored version in a new case, then monitor later runs for recurring failures.
\end{itemize}
\end{frame}

% Page 58
\begin{frame}{Review the Completed System\handsup}
\begin{center}
\resizebox{.94\textwidth}{!}{\begin{tikzpicture}[box/.style={draw=blue,fill=blue!4,rounded corners=2pt,text width=5.5cm,minimum height=1.25cm,align=center,font=\small,inner sep=6pt}]
\node[box] at (3,0) {Evidence\\trace and reproduce};
\node[box] at (10,0) {Control\\block or explain};
\node[box,draw=forestgreen,fill=forestgreen!5] at (3,-2) {Continuity\\resume verified work};
\node[box,draw=forestgreen,fill=forestgreen!5] at (10,-2) {Improvement\\evaluate and restore};
\end{tikzpicture}
}
\end{center}
\begin{itemize}
\item Question: Choose one control and explain whose need it addresses.
\item Show what has been tested, what the results support, and what still needs to be built.
\end{itemize}
\end{frame}
